\documentclass[12pt]{article}
\usepackage[utf8]{inputenc}
\usepackage[T1]{fontenc}
\usepackage[margin=1in]{geometry}
\usepackage{setspace}
\usepackage{xcolor}
\usepackage[colorlinks=true, urlcolor=blue!60!black, linkcolor=blue!60!black]{hyperref}
\setlength{\parindent}{2em}
\setlength{\parskip}{1.25em}
\renewcommand{\baselinestretch}{1.0}
\begin{document}
\begin{center}
{\LARGE \textit{Jackson v. Metropolitan Edison Co.}, 419 U.S. 345 (1974)\par}
\vspace{0.8em}
{\large Case Brief\par}
\end{center}
\vspace{1em}

\section{Memory Jogger}
A privately owned, heavily regulated electric monopoly that cut off a customer's power without notice or a hearing was not a state actor: the Court required a close nexus between the State and the specific challenged act, and confined the public-function doctrine to powers traditionally exclusively reserved to the State (\href{https://supreme.justia.com/cases/federal/us/419/345/}{Justia}).

\section{Detailed Case Facts}
Respondent Metropolitan Edison Co.\ (``Metropolitan'') was a privately owned and operated Pennsylvania corporation holding a certificate of public convenience issued by the Pennsylvania Public Utility Commission (``PUC''), empowering it to deliver electricity to a service territory that included the city of York, Pennsylvania. As a condition of holding its certificate, Metropolitan was subject to extensive state regulation governing rates, service standards, accounting, and facilities. Under Rule~15 of its general tariff filed with the PUC (Electric Pa.\ P.U.C.\ No.~41), Metropolitan reserved the right to discontinue electric service to any customer on reasonable notice for nonpayment of bills (and without notice for abuse, fraud, or meter tampering, though Metropolitan did not rely on those latter grounds in the lower courts).

Petitioner Catherine Jackson, a resident of York, received electric service at her home under an account in her own name until September 1970, when Metropolitan terminated the account for asserted delinquency in payments. A new account was then opened in the name of James Dodson, another occupant of the residence, and electric service resumed. The parties disputed whether payments due under the Dodson account were made. In August 1971, Dodson moved out of the residence; service continued thereafter, though concededly no payments were made and Jackson stated she received no bills during that period.

On October 6, 1971, Metropolitan employees visited the home to inquire about Dodson's whereabouts, which Jackson stated were unknown to her. The following day, October 7, 1971, another Metropolitan employee visited and told Jackson that the electric meter had been tampered with so that it failed to register electricity usage. Jackson disclaimed any knowledge of tampering and requested that the account be transferred from Dodson's name to that of ``Robert Jackson,'' later identified as her 12-year-old son. (As Justice Brennan noted in dissent, Jackson also offered to pay \$30 toward her delinquent 1970 bill to keep service running, which Metropolitan rejected.) Four days later, on October 11, 1971, without further notice to Jackson, Metropolitan employees disconnected her electric service.

\section{Procedural History}
Jackson sued Metropolitan in the United States District Court for the Middle District of Pennsylvania under 42 U.S.C.\ \S~1983, seeking damages and an injunction compelling Metropolitan to continue providing electricity until she received notice, a hearing, and an opportunity to pay any amounts found due. She argued that Pa.\ Stat.\ Ann., Tit.~66, \S~1171 (1959) conferred a state-law entitlement to reasonably continuous electric service and that Metropolitan's termination under its PUC-filed tariff constituted ``state action'' depriving her of property without procedural due process under the Fourteenth Amendment.

The District Court granted Metropolitan's motion to dismiss for failure to establish state action (348 F.\ Supp.\ 954 (M.D.\ Pa.\ 1972)), observing that the PUC's sole connection to the termination was Metropolitan's notice filing and the Commission's inaction. The United States Court of Appeals for the Third Circuit affirmed (483 F.2d 754 (3d Cir.\ 1973)). Noting a circuit split over whether utility shutoffs by regulated private companies constitute state action, the Supreme Court granted certiorari (415 U.S. 912 (1974)).

\section{Judicial Votes}
The Supreme Court \textbf{affirmed} the Third Circuit by a \textbf{6--3} vote. The case was argued on October 15, 1974, and decided on December 23, 1974 (\href{https://supreme.justia.com/cases/federal/us/419/345/}{Justia}):

\begin{itemize}
    \item \textbf{Majority Opinion:} Justice Rehnquist delivered the opinion of the Court, joined by Chief Justice Burger and Justices Stewart, White, Blackmun, and Powell.
    \item \textbf{Concurring Opinions:} None.
    \item \textbf{Dissenting Opinions:} Justice Douglas, Justice Brennan, and Justice Marshall each filed a dissenting opinion.
\end{itemize}

\section{Holding}
\textbf{Questions Presented:}

Whether the termination of residential electric service for alleged nonpayment by a privately owned, heavily regulated public utility---operating under a state-issued certificate of public convenience, enjoying at least a partial monopoly, and acting pursuant to a termination provision in a general tariff filed with the state public utility commission---constitutes ``state action'' attributable to the State for purposes of the Fourteenth Amendment's Due Process Clause and 42 U.S.C.\ \S~1983.

\textbf{Holding:} No. Affirming the Third Circuit in a 6--3 decision, the Court held that Pennsylvania was not sufficiently connected with Metropolitan's termination of Jackson's electric service to render the utility's conduct attributable to the State under the Fourteenth Amendment. To establish state action by a regulated private entity, a plaintiff must show a ``sufficiently close nexus between the State and the challenged action of the regulated entity so that the action of the latter may be fairly treated as that of the State itself'' (\href{https://supreme.justia.com/cases/federal/us/419/345/}{419 U.S. at 351}). Moreover, the ``public function'' doctrine applies only to the exercise by a private entity of powers ``traditionally exclusively reserved to the State'' (\href{https://supreme.justia.com/cases/federal/us/419/345/}{419 U.S. at 352}), which does not encompass furnishing utility services. Because state action was absent, the Court did not decide whether continued electric service was a protected ``property'' interest or what procedural safeguards due process would require if a State terminated service.

\section{Analysis of Opinions}
\subsection{Majority Opinion (Justice Rehnquist)}

Writing for the six-Justice majority, Justice Rehnquist began with the foundational dichotomy established in the \href{https://supreme.justia.com/cases/federal/us/109/3/}{\emph{Civil Rights Cases}}, 109 U.S. 3 (1883): the Fourteenth Amendment shields individuals against state deprivations, but ``offers no shield'' against private conduct, ``however discriminatory or wrongful'' (419 U.S. at 349). Rehnquist framed the dispute as whether Metropolitan's shutoff fell on the \href{https://supreme.justia.com/cases/federal/us/365/715/}{\emph{Burton v. Wilmington Parking Authority}}, 365 U.S. 715 (1961), side of the line (state action) or the \href{https://supreme.justia.com/cases/federal/us/407/163/}{\emph{Moose Lodge No. 107 v. Irvis}}, 407 U.S. 163 (1972), side (private action). Evaluating and rejecting each of Jackson's four theories in turn, the Court concluded that the case belonged with \emph{Moose Lodge}:

\subsection{Extensive State Regulation Alone Does Not Create State Action}
The mere fact that a private business is subject to ``extensive and detailed'' state regulation does not convert its conduct into state action (419 U.S. at 350, citing \href{https://supreme.justia.com/cases/federal/us/407/163/}{\emph{Moose Lodge}}, 407 U.S. at 176--177, and \href{https://supreme.justia.com/cases/federal/us/343/451/}{\emph{Public Utilities Comm'n v. Pollak}}, 343 U.S. 451, 462 (1952)). In Footnote~7, Justice Rehnquist underscored the overbreadth of a regulation-equals-state-action rule by noting that Pennsylvania's public-utility statute covered dozens of enterprises, including individually owned taxicabs; treating regulation alone as sufficient would mean ``the actions of a lone Philadelphia cab driver could also be fairly treated as those of the State of Pennsylvania'' (419 U.S. at 350 n.7). Rather than focusing on the general regulatory relationship, the constitutional inquiry demands a ``sufficiently close nexus'' between the State and the \emph{specific challenged action} (419 U.S. at 351).

\subsection{Monopoly Status Lacks a Nexus to the Challenged Termination}
Jackson contended that Pennsylvania had conferred monopoly status upon Metropolitan. The Court doubted this as a factual matter, explaining in Footnote~8 that Metropolitan was organized in 1874---39 years before Pennsylvania enacted utility regulation in 1913---and that utilities are ``natural monopolies'' created by high capital requirements and economies of scale, upon which regulation is superimposed as a substitute for competition rather than to eliminate it (419 U.S. at 351--352 n.8). Even assuming Pennsylvania had granted a monopoly, the Court held that monopoly status was not determinative because there was no sufficient relationship between Metropolitan's monopoly status and its specific decision to terminate Jackson's service, just as monopoly aspects had not established state action in \href{https://supreme.justia.com/cases/federal/us/343/451/}{\emph{Pollak}} or \href{https://supreme.justia.com/cases/federal/us/407/163/}{\emph{Moose Lodge}} (419 U.S. at 351--352).

\subsection{Narrowing the Public-Function Doctrine}
Jackson argued that supplying an essential public service mandated to be ``reasonably continuous'' under Pa.\ Stat.\ Ann., Tit.~66, \S~1171 constituted a ``public function.'' In the opinion's most consequential doctrinal move, Justice Rehnquist sharply narrowed the public-function exception to ``the exercise by a private entity of powers traditionally exclusively reserved to the State''---such as conducting elections (\emph{Nixon v. Condon}; \emph{Terry v. Adams}), operating a company town (\emph{Marsh v. Alabama}), or running a municipal park (\emph{Evans v. Newton})---or exercising delegated sovereign powers ``such as eminent domain'' (419 U.S. at 352--353). Because Pennsylvania law imposed no obligation on the State itself to furnish electricity and Pennsylvania courts had long held that utility service was neither a state function nor a municipal duty, supplying electricity was not ``traditionally the exclusive prerogative of the State'' (419 U.S. at 353).

The Court also rejected Jackson's fallback proposal to treat all businesses ``affected with the public interest'' as state actors. Quoting \emph{Nebbia v. New York}, 291 U.S. 502, 536 (1934), Rehnquist observed that ``affected with a public interest'' defines no closed category and encompasses ``[d]octors, optometrists, lawyers, Metropolitan, and \emph{Nebbia}'s upstate New York grocery selling a quart of milk'' (419 U.S. at 353--354). Providing essential, regulated goods or services does not transform private actors into agents of the State.

\subsection{Passive Tariff Filing Is Not State Authorization or Compulsion}
Jackson argued that Pennsylvania specifically authorized and approved Metropolitan's shutoff practice through Rule~15 of its filed general tariff. The Court found no state ``imprimatur'' (419 U.S. at 357). Rule~15 had appeared in Metropolitan's tariffs for years, reflected a pre-existing common-law right, and had never been the subject of a PUC hearing or scrutiny; although the Commission held hearings on a general rate increase in Tariff No.~41, Rule~15 simply took effect automatically 60 days after filing when the PUC did not disapprove it (419 U.S. at 354--355 \& nn.10--13).

The Court distinguished \href{https://supreme.justia.com/cases/federal/us/343/451/}{\emph{Public Utilities Comm'n v. Pollak}}, 343 U.S. 451 (1952), noting first that \emph{Pollak} had arguably only assumed state action \emph{arguendo}, and second that the D.C.\ utility commission in \emph{Pollak} had affirmatively investigated piped radio music on transit buses on its own motion, held a full hearing, and concluded that the practice positively improved passenger conditions (419 U.S. at 356--357). By contrast, regulatory approval or acquiescence in a utility-initiated practice---``where the commission has not put its own weight on the side of the proposed practice by ordering it''---does not convert a private business choice into state action (419 U.S. at 357).

\subsection{Absence of a Symbiotic Relationship}
Finally, the Court distinguished \href{https://supreme.justia.com/cases/federal/us/365/715/}{\emph{Burton v. Wilmington Parking Authority}}, 365 U.S. 715 (1961), limiting \emph{Burton}'s ``symbiotic relationship'' holding to lessees of public property (419 U.S. at 357--358). Unlike the Eagle Coffee Shoppe operating inside a publicly owned parking garage, Metropolitan was a privately owned corporation that did not lease facilities from Pennsylvania, bore sole responsibility for generating and delivering power, and merely paid taxes and complied with regulation like other corporations. Viewing all of Jackson's arguments together, the Court concluded that Pennsylvania was not a ``partner or even a joint venturer'' in Metropolitan's termination decision (419 U.S. at 358, quoting \href{https://supreme.justia.com/cases/federal/us/407/163/}{\emph{Moose Lodge}}, 407 U.S. at 176--177).

\subsection{Dissenting Opinion (Justice Douglas)}

Justice Douglas attacked the majority's methodology as ``fundamentally sequential, rather than cumulative'' (419 U.S. at 362--363). Invoking \href{https://supreme.justia.com/cases/federal/us/365/715/}{\emph{Burton}}'s command to sift facts and weigh circumstances, Douglas argued that courts must evaluate the \emph{aggregate} of state-private ties rather than isolating and dismissing each factor seriatim. Taken together, Pennsylvania's licensing of a sole-provider monopolist supplying a modern necessity of life (without which a home ``is likely to become unlivable''), its comprehensive regulatory control, and its statutory empowerment of the utility's tariff rules made Metropolitan's conduct state action under 42 U.S.C.\ \S~1983 (419 U.S. at 361--363). Drawing on Justice Brandeis and a 1936 Federal Trade Commission report, Douglas warned that state utility commissions had largely become ``prisoners of the utilities'' (419 U.S. at 363) and highlighted the irony that utilities routinely invoked ``state action'' immunity to shield themselves from federal antitrust liability while disclaiming state action when sued under \S~1983.

\subsection{Dissenting Opinion (Justice Brennan)}

Justice Brennan declined to reach the state-action question because he believed no justiciable controversy existed. Under \S~25 of Metropolitan's tariff, Metropolitan's statutory duty under \S~1171 ran only to ``customers'' lawfully receiving service. Because Jackson's own account had been terminated in September 1970 and James Dodson was the customer of record between 1970 and 1971, Brennan concluded Jackson lacked any state-law entitlement to service when Metropolitan rejected her \$30 partial-payment offer and shut off power on October 11, 1971. He would have vacated the judgment and remanded for dismissal without opining on state action. (In Footnote~2, the majority responded that whether Jackson held a protected ``property'' entitlement under \emph{Board of Regents v. Roth} was itself a merits constitutional question no less than state action.)

\subsection{Dissenting Opinion (Justice Marshall)}

Justice Marshall agreed with Justice Brennan that the confusing factual record made the case a poor vehicle and would have dismissed the writ as improvidently granted. Because the majority reached the merits, however, Marshall wrote a forceful substantive dissent resting on three core arguments:
\begin{enumerate}
  \item \textbf{State-Sanctioned Monopoly and Interdependence:} Marshall rejected the majority's ``natural monopoly'' distinction as speculative economics, emphasizing that Pennsylvania deliberately foreclosed competition and partnered with Metropolitan so its service would be ``the functional equivalent of service provided by the State'' (419 U.S. at 367--368), including delegating eminent domain and property-entry rights (419 U.S. at 368 n.3/1).
  \item \textbf{Rejection of an ``Ordering'' Requirement:} Requiring the PUC to \emph{order} a practice rather than authorize or approve it broke sharply with precedent from the \href{https://supreme.justia.com/cases/federal/us/109/3/}{\emph{Civil Rights Cases}} through \href{https://supreme.justia.com/cases/federal/us/387/369/}{\emph{Reitman v. Mulkey}}, 387 U.S. 369 (1967), and effectively gutted \href{https://supreme.justia.com/cases/federal/us/343/451/}{\emph{Pollak}} (419 U.S. at 369--371). It also trapped plaintiffs in a catch-22: approval without a hearing showed no state weight, while a commission challenge showed state disapproval.
  \item \textbf{Public Function, Pluralism, and Nondiscrimination:} Unlike parochial schools or private clubs---where immunity from constitutional rules protects private pluralism, diversity, and individual choice---``[t]he values of pluralism and diversity are simply not relevant when the private company is the only electric company in town'' (419 U.S. at 372--373). Marshall argued that courts could avoid burdensome hearings simply by requiring modest pre-termination notice and access to a responsible official (419 U.S. at 373). Finally, he warned that because the majority drew no distinction among constitutional rights, its holding would logically immunize a private electric monopoly that refused service on the basis of race or welfare status (419 U.S. at 373--374).
\end{enumerate}

\subsection{Subsequent Developments}

\begin{itemize}

\item \textbf{Consolidation (1978--1999):} The Court extended \emph{Jackson}'s exclusivity and compulsion limits across private economic and social-welfare disputes in \href{https://supreme.justia.com/cases/federal/us/436/149/}{\emph{Flagg Bros., Inc. v. Brooks}}, 436 U.S. 149 (1978) (warehouseman's UCC lien sale), \href{https://supreme.justia.com/cases/federal/us/457/991/}{\emph{Blum v. Yaretsky}}, 457 U.S. 991 (1982) (Medicaid-funded nursing homes), \href{https://supreme.justia.com/cases/federal/us/457/830/}{\emph{Rendell-Baker v. Kohn}}, 457 U.S. 830 (1982) (publicly funded private school), \href{https://www.law.cornell.edu/supremecourt/text/457/922}{\emph{Lugar v. Edmondson Oil Co.}}, 457 U.S. 922 (1982) (two-part fair-attribution test), and \href{https://www.law.cornell.edu/supremecourt/text/526/40}{\emph{American Manufacturers Mutual Insurance Co. v. Sullivan}}, 526 U.S. 40 (1999) (workers'-compensation utilization review).

\item \textbf{\href{https://www.law.cornell.edu/supremecourt/text/17-1702}{\emph{Manhattan Community Access Corp. v. Halleck}}, 587 U.S. 802 (2019):} Relying directly on \emph{Jackson} (419 U.S. at 350, 352--354), the Court held 5--4 that operating public-access cable channels is not a traditional, exclusive public function, and that municipal designation, heavy regulation, and rights-of-way access do not make a private speech forum a state actor.

\end{itemize}

\section{Examples: Future Applications}
\begin{enumerate}

\item \textbf{Same-Side (No State Action): A Medicaid-Funded Nursing Home Transfers a Resident.} A privately owned nursing home, licensed and closely regulated by the State and receiving most of its revenue from Medicaid, decides through its own physicians that a resident needs only a lower level of care and transfers her to another facility without a hearing. The State then reduces her Medicaid benefits to match the new placement. Under \textit{Jackson}, licensing, heavy regulation, and public funding do not make the home's decision the State's; the question is whether the State compelled or put its own weight behind the specific transfer, and here private medical judgment made it. \textit{Blum v. Yaretsky}, 457 U.S. 991 (1982), applied \textit{Jackson} to these facts and found no state action, so the resident has no Fourteenth Amendment due process claim against the home.

\item \textbf{Same-Side (No State Action): A Franchised Internet Provider Cancels Service.} A privately owned broadband company holds an exclusive municipal franchise, strings its fiber along public rights-of-way, and is regulated by the state utility commission. It terminates a subscriber's service for disputed late payments without notice, under a termination clause in its filed terms of service that the commission never reviewed. Each of the subscriber's arguments fails as it did in \textit{Jackson}: regulation and a franchise do not supply a nexus to the specific termination; supplying internet service, like supplying electricity, has never been a power traditionally and exclusively reserved to the State; and the commission's passive acceptance of the filed terms does not put the State's weight behind the shutoff. Without state action, the subscriber's remedies lie in contract and state utility law.

\item \textbf{Opposite-Side (State Action): A Commission-Ordered Shutoff.} A state utility commission issues an emergency order directing private electric utilities to remotely disconnect high-usage smart meters without a hearing. Unlike \emph{Jackson}'s passive tariff filing, state compulsion establishes a close nexus (419 U.S. at 357; \href{https://supreme.justia.com/cases/federal/us/457/991/}{\emph{Blum}}), triggering \href{https://supreme.justia.com/cases/federal/us/436/1/}{\emph{Craft}} due-process scrutiny.

\item \textbf{Opposite-Side (Likely State Action): A Utility Condemns an Easement.} The same kind of private electric company, using a power of eminent domain that state law delegates to certificated utilities, files a condemnation proceeding and takes a permanent easement across a homeowner's backyard for a new transmission line. The homeowner sues, claiming the taking was not for public use and that the procedures denied her due process. \textit{Jackson} itself marks this as the contrasting case: the Court said that if Metropolitan were exercising ``some power delegated to it by the State which is traditionally associated with sovereignty, such as eminent domain,'' the case ``would be quite a different one.'' Condemnation is a sovereign power the State has lent to the utility, so the utility's use of it is fairly attributable to the State.

\item \textbf{Fence-Sitter (Charter School Unresolved): Charter School vs.\ Prison Clinic.} A 100\% state-funded private charter school expels a student without a hearing; a private prison healthcare contractor withholds insulin. The prison contractor is a state actor under \href{https://supreme.justia.com/cases/federal/us/487/42/}{\emph{West v. Atkins}}; the charter school sits on the unresolved \emph{Jackson}/\href{https://supreme.justia.com/cases/federal/us/457/830/}{\emph{Rendell-Baker}}/\href{https://www.law.cornell.edu/supct/cert/24-394}{\emph{Drummond}} boundary between non-exclusive education and statutory public-entity entwinement (\href{https://www.law.cornell.edu/supremecourt/text/531/288}{\emph{Brentwood}}). The school is the hard case. On one side, education has never been an exclusive state prerogative, and \textit{Rendell-Baker} held that a privately run school funded almost entirely by public money was not a state actor when it fired employees. On the other, a State that labels charter schools as public schools makes them part of its public education system, and controls their charters looks more like the entwined State in \textit{Brentwood}. Because the Court divided evenly in \textit{Drummond}, the outcome is open.

\end{enumerate}

\section{Critique}
\textbf{Core Critiques:} Critics emphasize three tensions: (1) \textbf{seriatim disaggregation} (Douglas, J.), dismissing regulation, monopoly, and tariff approval individually rather than cumulatively; (2) the \textbf{legal-realist baseline problem}, treating utility self-help and PUC silence as ``neutral'' private action despite positive state-law backing; and (3) \textbf{weak pluralism justification} (Marshall, J.), shielding an investor-owned monopoly rather than individual associational liberty.

\textbf{The Formalist Cliff (\textit{Memphis Light} Contrast):} Four years after \emph{Jackson}, \href{https://supreme.justia.com/cases/federal/us/436/1/}{\emph{Memphis Light, Gas \& Water Div. v. Craft}}, 436 U.S. 1 (1978), reached the opposite result on nearly identical utility-shutoff facts solely because the utility was \textbf{municipally owned}. With state action automatic, the Court held that Tennessee's ``for cause'' rule created a protected property interest requiring pre-termination notice and an opportunity to contest billing errors before designated personnel. The pairing highlights \emph{Jackson}'s formalist cliff: identical shutoffs by sole-provider monopolies trigger Fourteenth Amendment due process or none at all depending entirely on municipal versus private corporate ownership.

\section{Key Quotations}
All quotations below are verified verbatim against the primary opinion text at \href{https://supreme.justia.com/cases/federal/us/419/345/}{Justia (419 U.S. 345)}.

\begin{quote}
``But the inquiry must be whether there is a sufficiently close nexus between the State and the challenged action of the regulated entity so that the action of the latter may be fairly treated as that of the State itself.''
\end{quote}
\noindent --- Justice Rehnquist, for the Court, \href{https://supreme.justia.com/cases/federal/us/419/345/}{419 U.S. at 351}.

\begin{quote}
``We have, of course, found state action present in the exercise by a private entity of powers traditionally exclusively reserved to the State. \dots\ If we were dealing with the exercise by Metropolitan of some power delegated to it by the State which is traditionally associated with sovereignty, such as eminent domain, our case would be quite a different one.''
\end{quote}
\noindent --- Justice Rehnquist, for the Court, \href{https://supreme.justia.com/cases/federal/us/419/345/}{419 U.S. at 352--353}.

\begin{quote}
``Approval by a state utility commission of such a request from a regulated utility, where the commission has not put its own weight on the side of the proposed practice by ordering it, does not transmute a practice initiated by the utility and approved by the commission into `state action.'\,''
\end{quote}
\noindent --- Justice Rehnquist, for the Court, \href{https://supreme.justia.com/cases/federal/us/419/345/}{419 U.S. at 357}.

\begin{quote}
``It is not enough to examine seriatim each of the factors upon which a claimant relies and to dismiss each individually as being insufficient to support a finding of state action. It is the aggregate that is controlling.''
\end{quote}
\noindent --- Justice Douglas, dissenting, \href{https://supreme.justia.com/cases/federal/us/419/345/}{419 U.S. at 360}.

\begin{quote}
``The values of pluralism and diversity are simply not relevant when the private company is the only electric company in town. \dots\ Thus, the majority's analysis would seemingly apply as well to a company that refused to extend service to Negroes, welfare recipients, or any other group that the company preferred, for its own reasons, not to serve.''
\end{quote}
\noindent --- Justice Marshall, dissenting, \href{https://supreme.justia.com/cases/federal/us/419/345/}{419 U.S. at 372--373, 374}.


\end{document}
